\documentclass[11pt,a4paper]{article}
\usepackage[T1]{fontenc}
\usepackage[utf8]{inputenc}
\usepackage{lmodern}
\usepackage[margin=1in]{geometry}
\usepackage{amsmath,amssymb,amsthm}
\usepackage{microtype}
\usepackage{booktabs}
\usepackage{enumitem}
\usepackage{xurl}
\usepackage[colorlinks=true,linkcolor=blue,citecolor=blue,urlcolor=blue]{hyperref}
\newcommand{\SAT}{\mathsf{SAT}}
\newcommand{\PP}{\mathsf{P}}
\newcommand{\NP}{\mathsf{NP}}
\newcommand{\Ppoly}{\mathsf{P}/\mathsf{poly}}
\newcommand{\SIZE}{\mathsf{SIZE}}
\newcommand{\bits}{\{0,1\}}
\newcommand{\loc}{L^{U}}
\newcommand{\poly}{\operatorname{poly}}
\newcommand{\timeU}{\operatorname{time}_{U}}
\newcommand{\K}{K}
\newtheorem{theorem}{Theorem}[section]
\newtheorem{proposition}[theorem]{Proposition}
\newtheorem{lemma}[theorem]{Lemma}
\newtheorem{corollary}[theorem]{Corollary}
\theoremstyle{definition}
\newtheorem{definition}[theorem]{Definition}
\newtheorem{conjecture}[theorem]{Conjecture}
\theoremstyle{remark}
\newtheorem{remark}[theorem]{Remark}
\newcommand{\claimid}[1]{\textnormal{\textsf{[#1]}}}

\hypersetup{pdftitle={Local Descriptions of SAT Truth Tables},
pdfauthor={Matias Donoso Quiroga}}
\title{Local Descriptions of SAT Truth Tables:\\
A Framework for Time--Description Trade-offs}
\author{Matías Donoso Quiroga}
\date{Working draft v0.1.0 --- September 29, 2026}
\begin{document}
\maketitle
\begin{abstract}
We formulate a research framework for representing SAT truth tables by short
descriptions that answer individual queries within a specified time bound.
For an efficiently universal deterministic evaluator, polynomial description
length and polynomial local evaluation time characterize polynomial advice,
and hence polynomial-size Boolean circuits. We give self-contained expository
proofs of this characterization, the corresponding lower-bound quantifiers,
and an elementary residual-function bound for ordered decision diagrams.
A unit-clause satisfiability example has exponential width under a grouped
input order and constant width under an interleaved order, illustrating why
representation restrictions matter. Deterministic finite experiments reproduce
the latter example. We distinguish these established observations from the
open objective of excluding all polynomial-size, polynomial-time local
descriptions of SAT. This draft claims neither a new general circuit lower
bound nor a resolution of \(P\) versus \(NP\).
\end{abstract}
\noindent\textbf{Keywords:} SAT; local decoding; nonuniform computation;
Kolmogorov complexity; circuit complexity; ordered decision diagrams.
\section{Scope and status}
The motivating question is whether the information in a complete table of SAT
answers can determine the least computational effort needed to answer a new
SAT query. We grant access to the complete finite table and do not charge for
discovering a representation. This convention makes the resulting model
nonuniform: different input lengths may receive unrelated descriptions.
The object of study is the feasible relationship between description
length and query time.

The initial version is an expository foundation and a research protocol.
Its proved statements are standard consequences of computability, advice
complexity, counting, and residual-function arguments. No originality is
claimed for them. Proofs expose model assumptions and quantifiers.
The proposed lower bound for unrestricted local descriptions is an open
target, not a result of the paper.

Three quantities remain separate: the length of a stored description,
the work space of an evaluator, and its running time. A short program can
describe a long computation, while a computation can reuse its work space.
Exact deterministic bit recovery is a particular notion of prediction;
it is not a claim about probabilistic prediction, average log-loss, or
Shannon entropy. No unrestricted equivalence between those statistical
quantities and worst-case SAT running time is assumed.

Every substantive statement has a stable identifier, recorded in the
repository claim register. The register distinguishes expository proofs,
literature results with hypotheses, open targets, and finite observations.
Independent human review and proof-assistant verification are not claimed.

\section{A local-description model}
\subsection{Inputs, tables, and evaluators}
Fix a conventional explicit binary encoding of Boolean formulas.
Let \(\SAT(x)=1\) if \(x\) is a valid encoding of a satisfiable formula,
and let \(\SAT(x)=0\) otherwise. For \(n\geq1\), the string
\[
 T_n=(\SAT(x))_{x\in\bits^n}
\]
lists answers in lexicographic order. Its length is \(N=2^n\).
We use \(T_n[x]\) for the entry indexed by the \(n\)-bit string \(x\).
All polynomial query-time bounds refer to \(n\), not to \(N\).

Fix a deterministic, efficiently universal multitape evaluator \(U\) with
read-only inputs \(d,x\), ordinary work tapes, and a one-bit output.
Efficient universality means that a fixed algorithm with auxiliary data can
be encoded with constant program overhead and simulated with polynomial
overhead in its time and input lengths. All interpretation and access costs
are charged; \(U\) has no SAT oracle or uncharged computational primitive.
The simulation polynomial may depend on the simulated algorithm.

\begin{definition}[Local description length]\label{def:local}
For a total time bound \(t:\mathbb N\to\mathbb N\), set
\[
 \loc_t(T_n)=\min\left\{
 |d|:\ \forall x\in\bits^n,\quad
 U(d,x)=T_n[x]\ \text{and}\ \timeU(d,x)\leq t(n)
 \right\}.
\]
The minimum of an empty set is \(+\infty\).
The same description must answer every query of length \(n\).
\end{definition}

The description \(d_n\) may depend on \(n\). There is no requirement that
\(n\mapsto d_n\) be computable or efficiently constructible.
The description length \(|d_n|\) is not a definition of work-space usage.
A fresh description for each individual query is not allowed: it could
simply supply that query's answer and erase the computational question.

Local access is closely related to the established resource-bounded
Kolmogorov measure \(\mathrm{KT}\) and to circuit size
\cite{AllenderEtAl2006}. We retain two parameters rather than assert
an exact numerical identity with a particular convention for
\(\mathrm{KT}\). In particular, \(\mathrm{KT}\), Levin's \(\mathrm{Kt}\),
and the full-output measure \(K^t\) are different definitions.

\subsection{What unrestricted compression already gives}
Let \(K\) denote prefix-free Kolmogorov complexity for a fixed optimal machine.
\begin{proposition}[Short unrestricted descriptions \claimid{P01}]
\label{prop:short}
For the fixed SAT encoding,
\[
 K(T_n)\leq K(n)+O(1)=O(\log(n+2)).
\]
There is also a constant-length local description answering \(T_n[x]\)
for every \(n\), if no query-time bound is imposed.
\end{proposition}
\begin{proof}
A fixed procedure, supplied with \(n\), enumerates \(\bits^n\), rejects
malformed encodings, decides each remaining SAT instance by exhaustive
search, and writes the resulting table. A shortest self-delimiting
program for \(n\), followed by a fixed interpreter for this procedure,
proves the first inequality. The usual self-delimiting encoding of
an integer gives the asymptotic bound.
For local access, a fixed exhaustive-search SAT program takes \(x\)
directly as input. Its code does not depend on \(n\).
\end{proof}

This proposition holds whether or not \(\PP=\NP\). It is an
upper bound on description length, not an efficient decoding guarantee.
It also holds with SAT replaced by any fixed decidable language.
The unrestricted table complexity includes a description of the length;
the local evaluator can obtain the query length from its input.

\section{Polynomial descriptions characterize polynomial advice}
The class \(\Ppoly\) consists of languages decidable in polynomial time
with advice of polynomial length depending only on input length.
Equivalently, these are languages with polynomial-size Boolean circuit
families: hardwire advice into computation circuits, or evaluate an
encoded circuit supplied as advice \cite{AroraBarak2009}.

\begin{theorem}[Advice characterization \claimid{T01}]
\label{thm:advice}
Under the evaluator convention of Definition~\ref{def:local},
\[
 \SAT\in\Ppoly
 \quad\Longleftrightarrow\quad
 \exists a,b\in\mathbb N_{>0}\ \exists n_0\
 \forall n\geq n_0:\quad \loc_{n^b}(T_n)\leq n^a.
\]
\end{theorem}
\begin{proof}
Suppose first that the right-hand side holds. For each \(n\geq n_0\),
choose a witnessing description \(d_n\). Use \(d_n\) as advice and run
the fixed evaluator \(U(d_n,x)\). Both advice length and time are
polynomial in \(n\). The finitely many smaller input lengths can be
handled by a finite table in the deciding machine. Thus
\(\SAT\in\Ppoly\).

Conversely, suppose a fixed deterministic machine \(M\) decides
SAT using advice \(a_n\) of length at most \(q(n)\) in time at most
\(p(n)\), for polynomials \(p,q\). Encode \(M\), a self-delimiting
header, and \(a_n\) in \(d_n\). Its length is polynomial in \(n\).
Efficient universality supplies a polynomial time bound for
\(U(d_n,x)\). Increasing integer exponents and discarding finitely
many small \(n\) absorbs multiplicative constants, headers, and
simulation overhead.
\end{proof}

This is a polynomial-resource equivalence, not equality of program
length and gate count. The proof applies equally to any language.

\begin{corollary}[The exact lower-bound target \claimid{C01}]
\label{cor:target}
The following assertions are equivalent:
\begin{enumerate}[label=(\roman*)]
 \item \(\SAT\notin\Ppoly\).
 \item For every \(a,b\in\mathbb N_{>0}\) and every \(n_0\), there is
 \(n\geq n_0\) such that \(\loc_{n^b}(T_n)>n^a\).
\end{enumerate}
Either assertion implies \(\PP\ne\NP\).
\end{corollary}
\begin{proof}
Negating the quantified right-hand side of Theorem~\ref{thm:advice}
gives (ii). If \(\PP=\NP\), SAT has a polynomial-time algorithm and
therefore polynomial-size circuits, contradicting (i).
\end{proof}

Assertion (ii) quantifies over infinitely many input lengths, rather
than all sufficiently large lengths. An eventual lower bound would
also suffice but is not the logical negation above.
The converse implication from \(\PP\ne\NP\) to (i) is not established.
The target excludes nonuniform advice as well as uniform algorithms.

A single fixed description correctly answering all input lengths in
polynomial time is an ordinary polynomial-time SAT algorithm.
A new polynomial-size description at each length is extra freedom.
Ruling out only maximally compressed or logarithmic-length descriptions
would leave open longer polynomial descriptions with faster evaluation.

\section{What counting and structural identities do not establish}
\begin{lemma}[Counting local descriptions \claimid{L01}]
\label{lem:count}
Fix \(n,t\) and an integer \(s\geq0\). At most \(2^{s+1}-1\)
truth tables on \(n\) input bits have \(\loc_t(T)\leq s\).
For a uniformly random such table,
\[
 \Pr[\loc_t(T)\leq s]\leq\min\{1,2^{s+1-2^n}\}.
\]
\end{lemma}
\begin{proof}
There are \(\sum_{j=0}^s2^j=2^{s+1}-1\) descriptions of length at
most \(s\). A deterministic description satisfying the time bound
on every query determines at most one table. Divide the bound by
the total \(2^{2^n}\) possible tables.
\end{proof}
This statement identifies most tables as hard, not SAT's particular
table. Proposition~\ref{prop:short} explains why raw information content
cannot identify the desired computational obstruction. Shannon entropy
also requires a specified distribution; the size of one table is not
its computational hardness.

\begin{proposition}[Self-reduction requires boundary conditions \claimid{P02}]
\label{prop:boundary}
Let \(P\) be Boolean-valued on formulas, allowing Boolean constants.
Suppose for every variable \(z\) occurring in \(F\),
\[
 P(F)=P(F[z=0])\lor P(F[z=1]).
\]
If \(P\) equals direct evaluation on variable-free formulas, then
\(P(F)=\SAT(F)\) for every \(F\). Without that boundary condition,
both constant functions satisfy the recurrence.
\end{proposition}
\begin{proof}
Induct on the number of distinct variables. The boundary condition
proves the base case. Substitution removes \(z\), so induction
identifies both terms on the right as satisfiability answers; their
disjunction is satisfiability of \(F\). For a constant function with
value \(c\), the recurrence is \(c=c\lor c\).
\end{proof}
Uniqueness under identities and boundary conditions does not bound
circuit size. Substituted formulas also have different encoding
lengths: a circuit-family argument must manage these lengths or
supply explicit satisfiability-preserving padding.
Counting identities as independent gates is unjustified because
circuits can share subcomputations.

A deterministic halting computation using space \(s(n)\geq\log n\)
has exponentially many possible configurations in \(s(n)\), for a
fixed standard machine. Counting them gives an exponential
\emph{upper} bound on time. It does not force a computation to visit
all configurations and is not a time lower bound.

\section{A restricted model: ordered decision diagrams}
Ordered binary decision diagrams are classical representations
\cite{Bryant1986}. We use deterministic \emph{layered} diagrams:
each layer reads the next bit in a fixed order, every path reads
all bits, and the final state supplies the answer. Equivalent
residual computations may merge. This convention avoids ambiguity
about skipped levels and early terminals.

For \(a,b\in\bits^m\), define
\[
 g_m(a,b)=\bigwedge_{i=1}^m\neg(a_i\land b_i).
\]
This is the satisfiability function of the unit-clause formula
\[
 F_{a,b}=
 \left(\bigwedge_{i:a_i=1}z_i\right)\land
 \left(\bigwedge_{i:b_i=1}\neg z_i\right).
\]
The experiment uses the incidence encoding \((a,b)\), not the
general formula encoding of \(T_n\). Conversion to conventional
syntax does not by itself transfer an order-dependent bound to
arbitrary encodings or circuits.

\begin{lemma}[Residual-function width bound \claimid{L02}]
\label{lem:residual}
At a fixed input cut, let \(R\) be the number of distinct functions
of the unread suffix induced by assignments to the read prefix.
Every correct deterministic layered diagram has at least \(R\)
reachable states there. A diagram whose states are residual
functions attains the bound.
\end{lemma}
\begin{proof}
If two prefixes inducing different residual functions reached the
same state, a common suffix would require different answers. A
deterministic diagram cannot produce both. For attainability, take
one state for every residual function at every layer. Each edge
restricts the next variable, defining consistent transitions and
the correct terminal value.
\end{proof}

\begin{proposition}[Grouped-order lower bound \claimid{P03}]
\label{prop:grouped}
For order \(a_1,\ldots,a_m,b_1,\ldots,b_m\), the minimum width
after reading \(a\) is \(2^m\).
\end{proposition}
\begin{proof}
For distinct \(a,a'\), choose \(i\) with \(a_i\ne a'_i\).
Let \(e_i\) have only its \(i\)-th bit set. Then
\(g_m(a,e_i)=1-a_i\) and \(g_m(a',e_i)=1-a'_i\).
All \(2^m\) prefixes induce different residual functions.
Apply Lemma~\ref{lem:residual}.
\end{proof}

\begin{proposition}[Interleaved-order upper bound \claimid{P04}]
\label{prop:interleaved}
For order \(a_1,b_1,\ldots,a_m,b_m\), every layer needs at most
three states. Thus a layered diagram has \(O(m)\) nodes.
\end{proposition}
\begin{proof}
After a completed pair, a prefix either already contains a conflict
or leaves the same disjointness function on the remaining pairs.
There are at most two residual functions. After the first bit of
the next pair, the possibilities are a previous conflict, an
unrestricted next \(b_i\), or the requirement \(b_i=0\).
There are at most three residual functions.
Apply Lemma~\ref{lem:residual} to the \(2m+1\) layers.
\end{proof}

The function also has a linear-size general circuit. The exponential
grouped-order bound concerns a restricted representation. An
explicit diagram needs many nodes in that order; this says nothing
comparable about the length of a program generating the diagram.

\begin{table}[htbp]
\centering
% Generated by experiments/obdd_residuals.py; do not edit.
\begin{tabular}{rrrr}
\toprule
$m$ & Grouped cut & Grouped maximum & Interleaved maximum \\
\midrule
1 & 2 & 2 & 2 \\
2 & 4 & 4 & 3 \\
3 & 8 & 8 & 3 \\
4 & 16 & 16 & 3 \\
5 & 32 & 32 & 3 \\
6 & 64 & 64 & 3 \\
7 & 128 & 128 & 3 \\
\bottomrule
\end{tabular}

\caption{Exhaustive residual-function counts. Grouped cut is the width
after the first \(m\) bits; interleaved maximum includes all layers,
including terminals. Finite observations illustrate the proofs.}
\label{tab:widths}
\end{table}

\section{Related directions and the open target}
Connections between resource-bounded descriptions and lower bounds
must retain their models, time scales, promises, and hypotheses.

Chapman and Williams \cite{ChapmanWilliams2015} characterize
\(\NP\nsubseteq\Ppoly\) through natural properties designed against
SAT-solving circuits with specified checking behavior.
This concerns a specialized class of solvers, not arbitrary
properties separating SAT from arbitrary circuits. It is a
literature direction, not a lower-bound method supplied here.

Oliveira, Pich, and Santhanam \cite{OliveiraPichSanthanam2021},
Theorem~1.4, prove the following hardness-magnification implication.
There is a universal \(c\geq1\) such that, if some \(\varepsilon>0\)
satisfies, for every sufficiently small constant \(\beta>0\),
\[
 \mathsf{Gap\text{-}MCSP}
 \left[\frac{N^\beta}{c\log_2 N},N^\beta\right]
 \notin\SIZE[N^{1+\varepsilon}],
\]
then \(\NP\nsubseteq\Ppoly\).
Here \(N\) is the length of the entire input truth table. The
promise distinguishes circuit complexity at most the first threshold
from complexity at least the second.
A restricted diagram bound for \(T_n\) does not satisfy this
antecedent. An application requires a separate resource-preserving
reduction.

Kabanets and Kolokolova \cite{KabanetsKolokolova2026} characterize
specified chain rules for time-bounded Kolmogorov complexity.
The March 13, 2026 revision states a direct equivalence between a
conditional chain rule with sublinear error and \(\PP=\NP\),
\emph{assuming} NP-hardness of the stated
\(\mathsf{GapMcK}^{t}\mathsf P\) problem.
It also gives unconditional characterizations involving other
complexity assumptions. We neither remove the extra hypothesis nor
identify their full-output conditional measure with local access.

\begin{conjecture}[Research target \claimid{Q01}]
\label{conj:main}
For every \(a,b\in\mathbb N_{>0}\), there are infinitely many \(n\)
such that \(\loc_{n^b}(T_n)>n^a\).
\end{conjecture}
By Corollary~\ref{cor:target}, this is the established open target
\(\SAT\notin\Ppoly\) in local-description language, not a new
conjecture or a proof advance.

An intermediate result must name its restricted model and identify
the changes of representation it tolerates. Candidate directions
include order-independent residual bounds for specified diagram
classes, parameter-checked reductions to a meta-complexity promise
problem, and proof-assistant formalization of the equivalence.
These tasks remain open in this project.
A bound \(|d|t^\alpha\geq n^c\) for fixed \(c\), on its own, still
permits polynomial description length and polynomial time.

\section{Reproducibility and review}
The repository supplies dependency-free Python experiments, unit
tests, LaTeX sources, references, a claim register, and CI configuration.
The experiment constructs the truth table of \(g_m\), enumerates every
prefix and suffix in each selected order, and counts distinct residual
truth tables. The published range is \(1\leq m\leq7\).
All arithmetic is exact; there are no random seeds or statistical
estimates.

Tests compare the incidence predicate with an independent exhaustive
unit-clause SAT verifier, check the grouped and interleaved width
predictions on small inputs, and test a separate parity example.
The manuscript table is regenerated and checked for drift.
These checks validate finite artifacts, not the general mathematical
proofs.

The command \texttt{make check} runs code and research-metadata checks;
\texttt{make pdf} builds the paper with \texttt{latexmk};
\texttt{make verify} performs both and checks the LaTeX log.
The README lists the toolchain.
Independent review must inspect the evaluator model, uniformity
quantifiers, finite exceptions, and representation restrictions.
Journal or preprint submission remains an author-controlled future step.

\paragraph{Authorship and assistance.}
The research question was developed by Matías Donoso Quiroga.
OpenAI ChatGPT/Codex assisted in preparing this initial exposition,
code, and repository structure. Independent human mathematical
review is pending. AI systems are not listed as authors.
No institutional affiliation, publication venue, or endorsement
is asserted.

\bibliographystyle{plain}
\bibliography{references}
\end{document}
